\begin{codeblock}
namespace std {
  struct @\libglobal{wrapper}@ {
    // \ref{leading.members}, members
    wrapper() = default;

    virtual void run() = 0;
  };
}
\end{codeblock}

\rSec3[leading]{Leading descriptions}

\rSec4[leading.general]{General}

\pnum
A wrapper stores one value.

\pnum
Its description is routed independently of its synopsis.

\rSec4[leading.members]{Members}

\indexlibraryctor{wrapper}%
\begin{itemdecl}
wrapper() = default;
\end{itemdecl}

\begin{itemdescr}
\pnum
Constructs an empty wrapper.
\end{itemdescr}

\indexlibrarymember{run}{wrapper}%
\begin{itemdecl}
virtual void run() = 0;
\end{itemdecl}

\begin{itemdescr}
\pnum
A concrete wrapper supplies this operation.
\end{itemdescr}

\rSec4[leading.functions]{Functions}

\indexlibraryglobal{first}%
\begin{itemdecl}
inline void first();
\end{itemdecl}

\begin{itemdescr}
\pnum
The first function has prose only.
\end{itemdescr}

\indexlibraryglobal{second}%
\begin{itemdecl}
inline void second();
\end{itemdecl}

\begin{itemdescr}
\pnum
The second function has its own description.
\end{itemdescr}

\indexlibraryglobal{third}%
\begin{itemdecl}
inline void third();
inline void third(int);
\end{itemdecl}

\begin{itemdescr}
\pnum
The third function offers these choices:
\begin{itemize}
\item a first choice, or
\item a second choice.
\end{itemize}
It introduces the labelled requirements.

\pnum
This is a separate leading paragraph with \iref{leading.general}.

\pnum
\effects
Performs the operation.
\end{itemdescr}
